\chapter{模块与工程}
\label{ch:18}

前面十七章，程序都写在一个 \texttt{.\allowbreak{}mgl} 文件里。真实项目做不到——几千行塞一个文件，谁也没法读。这一章讲怎么把代码拆开，以及工具链给工程准备的那几个命令。

拆开有两条线：

\begin{itemize}

\item 语言层：\texttt{module} 声明一个文件属于哪个模块，\texttt{import} 把别的模块引进来；

\item 工程层：\texttt{magel new} 起一个带目录结构的工程，\texttt{magel test} / \texttt{magel fmt} 管测试和风格。

\end{itemize}

\section{\texttt{module}：给文件一个名字}

每个文件首行都可以声明自己属于哪个模块：

\par
\medskip
\begin{lstlisting}[style=magel]
module mylib::math;
\end{lstlisting}
\par
\medskip

\texttt{mylib:\allowbreak{}:\allowbreak{}math} 是模块路径，用 \texttt{:\allowbreak{}:\allowbreak{}} 分隔——和标准库里的 \texttt{std:\allowbreak{}:\allowbreak{}iter} 是同一套写法。

模块名不是随便起的：\texttt{import} 一个模块时，编译器要按这个名字去找对应的文件（见 18.2）。

核心模块 \texttt{std:\allowbreak{}:\allowbreak{}lang} 不用声明也不用 import：\texttt{println} / \texttt{format} / \texttt{assert} / \texttt{len} / \texttt{Any} / \texttt{Option} / \texttt{Result} 这些都是它自动带进来的（前面十七章你一直在用，从没 import 过）。

\section{\texttt{import}：把别的模块接进来}

假设有两个文件。先写一个工具模块 \texttt{util/\allowbreak{}text.\allowbreak{}mgl}：

\par
\medskip
\begin{lstlisting}[style=magel]
module util::text;

String shout(String s)
{
    return s + "!";
}
\end{lstlisting}
\par
\medskip

再用 \texttt{import} 在另一个文件里用它：

\par
\medskip
\begin{lstlisting}[style=magel]
module app;
import util::text;

int64 main()
{
    println("{}", shout("hi"));
    return 0;
}
\end{lstlisting}
\par
\medskip

跑 \texttt{magel run app.\allowbreak{}mgl}，输出：

\par
\medskip
\begin{lstlisting}[style=out]
hi!
\end{lstlisting}
\par
\medskip

几个要点：

\begin{itemize}

\item \texttt{import util:\allowbreak{}:\allowbreak{}text;} 之后，\texttt{util:\allowbreak{}:\allowbreak{}text} 里的公开函数按裸名就能调——直接写 \texttt{shout(\allowbreak{}"hi")\allowbreak{}}，不用写 \texttt{util:\allowbreak{}:\allowbreak{}text:\allowbreak{}:\allowbreak{}shout}；

\item 模块名 \texttt{util:\allowbreak{}:\allowbreak{}text} 对应文件 \texttt{util/\allowbreak{}text.\allowbreak{}mgl}：\texttt{:\allowbreak{}:\allowbreak{}} 换成目录分隔符，末尾加 \texttt{.\allowbreak{}mgl}。文件名和模块名对不上，\texttt{import} 就找不到；

\item 路径是相对当前文件解析的。

\end{itemize}

文件名 = 模块路径这条对应关系，是把「模块」和「目录结构」绑在一起的关键：看文件树就知道模块树。

\begin{tip}{提示}
\texttt{import} 后也可以写完整路径来调用，或者用 \texttt{import\_names} 引入短名——细节见语言参考第 12 节与 \texttt{example/\allowbreak{}samples/\allowbreak{}import\_names.\allowbreak{}mgl}。初学阶段先把 18.2 这套用熟。
\end{tip}

\section{起一个工程：\texttt{magel new}}

手写目录、手写 \texttt{magel.\allowbreak{}toml} 太麻烦。工具链给了一条命令：

\par
\medskip
\begin{lstlisting}[style=sh]
magel new demo
\end{lstlisting}
\par
\medskip

它创建一个叫 \texttt{demo} 的新目录（目标目录必须不存在，免得覆盖你的东西），里面长这样：

\par
\medskip
\begin{lstlisting}[style=out]
demo/
├── .gitignore
├── README.md
├── magel.toml          # 工程配置
├── src/
│   └── main.mgl        # 应用源码
└── tests/
    └── test_main.mgl   # 测试
\end{lstlisting}
\par
\medskip

建完还会顺手 \texttt{git init} 建好仓库。\texttt{src/\allowbreak{}main.\allowbreak{}mgl} 是那个你熟得不能再熟的 hello：

\par
\medskip
\begin{lstlisting}[style=magel]
int64 main()
{
    println("Hello, Magel!");
    return 0;
}
\end{lstlisting}
\par
\medskip

\section{\texttt{magel.\allowbreak{}toml}：工程配置文件}

工程靠一个 \texttt{magel.\allowbreak{}toml} 描述自己。默认生成的是：

\par
\medskip
\begin{lstlisting}[style=cfg]
[package]
name = "hello"
version = "0.1.0"

[magel]
run = "interpret"
main = "src/main.mgl"
\end{lstlisting}
\par
\medskip

\begin{itemize}

\item \texttt{[\allowbreak{}package]\allowbreak{}} 里是包名和版本；

\item \texttt{[\allowbreak{}magel]\allowbreak{}} 里告诉工具链：入口是哪个文件（\texttt{main}）、默认用什么方式跑（\texttt{run}）。

\end{itemize}

工程根目录一旦有 \texttt{magel.\allowbreak{}toml}，在这个目录里执行命令就不用反复写路径了。

\texttt{.\allowbreak{}gitignore} 也一并给好了——构建产物不该进版本库：

\par
\medskip
\begin{lstlisting}[style=out]
/build/
/.magel/
/magel_out*
*.exe
\end{lstlisting}
\par
\medskip

\section{三个日常命令}

工程起来了，日常就三件事：跑、查、测。

\par
\medskip
\begin{lstlisting}[style=sh]
magel run src/main.mgl     # 运行
\end{lstlisting}
\par
\medskip

\par
\medskip
\begin{lstlisting}[style=out]
Hello, Magel!
\end{lstlisting}
\par
\medskip

\texttt{run} 是解释执行（第 1 章讲过 \texttt{run} 和 \texttt{check} 的区别）：

\par
\medskip
\begin{lstlisting}[style=sh]
magel check src/main.mgl   # 只检查，不运行
\end{lstlisting}
\par
\medskip

\texttt{check} 只做编译期检查，一行都不执行——改完代码先 \texttt{check} 一下，比 \texttt{run} 快，也安全（不会真的动到数据）。

测试用 \texttt{test}。工程里的 \texttt{tests/\allowbreak{}test\_main.\allowbreak{}mgl} 是个普通的 \texttt{main}，里面全是 \texttt{assert} （第 16 章）：

\par
\medskip
\begin{lstlisting}[style=magel]
int64 main()
{
    assert(1 + 1 == 2);
    return 0;
}
\end{lstlisting}
\par
\medskip

跑整个 \texttt{tests/\allowbreak{}} 目录：

\par
\medskip
\begin{lstlisting}[style=sh]
magel test tests
\end{lstlisting}
\par
\medskip

输出每个测试文件的结果和合计（路径按实际位置显示）：

\par
\medskip
\begin{lstlisting}[style=out]
[passed] tests/test_main.mgl
通过 1，失败 0，跳过 0
\end{lstlisting}
\par
\medskip

\texttt{assert} 就是测试的体检项：条件不成立，那个测试文件就判 \texttt{失败}。你不用装任何测试框架——\texttt{main} + \texttt{assert} 就是最小可用的测试。

\section{\texttt{magel fmt}：让风格统一}

一个团队里，最没意义的争论是「大括号换行吗」「缩进几格」。工具链直接给了一个格式化器：

\par
\medskip
\begin{lstlisting}[style=sh]
magel fmt src/main.mgl            # 打印格式化后的源码到屏幕
magel fmt src/main.mgl --write    # 直接写回文件
magel fmt src/main.mgl --check    # 只检查「格式对不对」，不修改
\end{lstlisting}
\par
\medskip

\texttt{--check} 是给 CI 用的：文件已经规范 \ensuremath{\rightarrow} 安静退出；需要格式化 \ensuremath{\rightarrow} 打印文件名并以非零码退出：

\par
\medskip
\begin{lstlisting}[style=sh]
magel fmt src/messy.mgl --check
\end{lstlisting}
\par
\medskip

\par
\medskip
\begin{lstlisting}[style=out]
需要格式化：src/messy.mgl
\end{lstlisting}
\par
\medskip

（退出码是 1。CI 里挂上这条，谁提交了没格式化的代码就会红。）

\texttt{fmt} 会干的事，拿一段挤成一团的代码试（它本身能跑，只是很难看）：

\par
\medskip
\begin{lstlisting}[style=magel]
int64 main()
{
    int64 x = 1;
    if (x>0)
    {
        println("{}", x);
    }
    return 0;
}
\end{lstlisting}
\par
\medskip

\texttt{magel fmt} 之后：

\par
\medskip
\begin{lstlisting}[style=magel]
int64 main()
{
    int64 x = 1;
    if (x>0)
    {
        println("{}", x);
    }
    return 0;
}
\end{lstlisting}
\par
\medskip

缩进统一成 4 空格、层级对齐好。注意它只管空白和换行，不动代码内容 （\texttt{x=\allowbreak{}1} 里没有空格、\texttt{"\{\}",\allowbreak{}    x} 里那串多余空格，都原样留着）——格式化器不替你做美学判断，只保证结构清楚。

\subsection{终端查手册}

\texttt{magel man} 是终端手册入口，从同一可执行文件中的标准库数据查询。无目标时按能力域列模块，模块目标列符号，类型和方法目标显示签名、参数、返回值、错误边界及已有完整示例。

\par
\medskip
\begin{lstlisting}[style=sh]
magel man
magel man std::text
magel man std::text::Text::split
magel man --search delimiter
magel man --json std::text::Text::split
magel man std::text::Text::split --run
\end{lstlisting}
\par
\medskip

搜索默认忽略大小写，\texttt{--exact} 匹配完整字段；拼错目标会给候选。\texttt{--lang zh|en} 选择界面和诊断语言，库说明仍复用中文原文。\texttt{--run} 在临时目录运行完整示例并附上输出。普通查询不需要 Python、网络或仓库；开发用 \texttt{--update} 运行既有生成脚本，重新构建后才进入内嵌数据。

完整命令和 JSON 结构见 \href{https://github.com/rabbitc999/Magel/blob/master/doc/man/README.md}{doc/man}，标准库 API 说明仍在 \href{https://github.com/rabbitc999/Magel/blob/master/doc/library/README.md}{doc/library}。

\subsection{编辑器增量提示}

LSP 将补全、悬停、语义着色、工作区查询、宏预览、编译诊断交给不同工作队列。修改按 UTF-16 编辑范围同步，同一源码版本复用索引。诊断延迟 250 毫秒合并连续编辑，保存立即检查；取消过期任务，并且不把旧版本结果写回编辑器。

缓存在 magel 可执行文件旁的 \texttt{build/\allowbreak{}lsp/\allowbreak{}}，与工程目录无关。未保存的依赖修改和关闭会刷新引用者，不写回源文件。缓存自动清理七天未使用的普通文件，超过 256 MiB 时淘汰到 192 MiB，保护最近十分钟的文件；活跃任务可以暂时超过额度。索引内存也有淘汰阈值，缓存损坏或不可写会退回分析。详细策略见 \href{https://github.com/rabbitc999/Magel/blob/master/doc/man/lsp.md}{LSP 增量提示与缓存}。

\subsection{\texttt{inline}：函数定义的内联属性}

\texttt{inline} 写在模块级函数定义前，参数、返回类型与函数体仍按普通函数规则检查。

\par
\medskip
\begin{lstlisting}[style=magel]
inline int64 square(int64 value)
{
    return value * value;
}

int64 main()
{
    println("{}", square(3));
    return 0;
}
\end{lstlisting}
\par
\medskip

\par
\medskip
\begin{lstlisting}[style=out]
9
\end{lstlisting}
\par
\medskip

它是函数定义的属性，不是任意表达式的前缀，也不改变值语义。代码组织仍以清楚的职责为先，不需要给所有函数都加上它。

\section{小结}

\begin{center}
\begin{tabular}{ll}
\toprule
你会了 & 关键点 \\
\midrule
模块声明 & \texttt{module mylib:\allowbreak{}:\allowbreak{}math;} 放首行 \\
导入 & \texttt{import util:\allowbreak{}:\allowbreak{}text;}，公开函数按裸名调用 \\
名字对应文件 & \texttt{util:\allowbreak{}:\allowbreak{}text} \ensuremath{\leftrightarrow} \texttt{util/\allowbreak{}text.\allowbreak{}mgl}，路径相对当前文件 \\
\texttt{std:\allowbreak{}:\allowbreak{}lang} & 核心模块，自动加载，永不 import \\
起工程 & \texttt{magel new demo}（目标目录不能已存在） \\
工程配置 & \texttt{magel.\allowbreak{}toml} 写包名、入口、运行方式 \\
日常命令 & \texttt{run}（跑）、\texttt{check}（只查）、\texttt{test}（跑 \texttt{tests/\allowbreak{}}） \\
测试 & 普通 \texttt{main} + \texttt{assert}，无需框架 \\
格式化 & \texttt{magel fmt … [\allowbreak{}--check|--write]\allowbreak{}}，只动空白不碰代码 \\
\bottomrule
\end{tabular}
\end{center}

\section{下一步}

工程拆好了、测试挂上了、格式统一了。接下来两章讲工具链里更硬核的部分：接下来会见到 \texttt{\%define} / \texttt{\$macro} 这些带 \texttt{\%} 和 \texttt{\$} 的写法：它们是 Magel 的预处理与宏，能在编译前替你改源码。

翻到\hyperref[ch:19]{第 19 章 预处理与宏}。

书稿与关键字帮助也在程序内部：\texttt{magel man book} 导出书稿与图片到当前目录的 \texttt{book/\allowbreak{}} 并列出章节，\texttt{magel man book 6} 查看第 6 章，\texttt{magel man keyword break} 查看关键字语法与示例。\texttt{magel man skill} 输出 AI SKILL，可保存并交给支持 SKILL 的 AI 工具；VS Code 包也携带同一份 \texttt{skills/\allowbreak{}magel-programming/\allowbreak{}SKILL.\allowbreak{}md}。

\subsection{终端手册的入口}

\par
\medskip
\begin{lstlisting}[style=sh]
magel man api std
magel man api std::text::Text
magel man keyword return
magel man keyword '#derive'
magel man macro
\end{lstlisting}
\par
\medskip

API 查询必须显式写 \texttt{api}。省略它的 \texttt{magel man std} 或 \texttt{magel man std:\allowbreak{}:\allowbreak{}text:\allowbreak{}:\allowbreak{}Text} 会报错并给出正确命令；\texttt{keyword} 查询语言关键字，\texttt{macro} 查询宏规则，二者不是 API 入口。可以给这些查询加 \texttt{--json} 供工具读取。
